\scanpage{46}
% archive sheet number: 45
schémas abéliens, \inserted{permettant de définir un foncteur qu.inverse du foncteur du théorème.} Pour ceci, on considère pour tout schéma abélien $A$ sur $S$ son e xtension vectorielle universelle, qui est une extension
\[
0\handwritten{\longrightarrow}\underline{t}_{A^*}^{\handwritten{\vee}}\handwritten{\unsure{\mathbin{-}|\longrightarrow}}E(A)\handwritten{\mathbin{-}}A\handwritten{\longrightarrow}0\quad .
\]
Il se trouve qu'on peut trouver un cristal en groupes lisses canonique sur $S$, soit $\underline{E}(\deleted{\unsure{S}}A)$, tel que l'on ait
\[
E(A)\handwritten{\unsure{\simeq|\overset{\sim}{\longrightarrow}}}\underline{E}(A)_S\quad .
\]
La définition est telle que \deleted{pour tout le foncteur}\inserted{$A$} $\underline{E}(A)$ définisse un foncteur
\[
\text{Schémab}(S)\handwritten{\longrightarrow}\text{Cris Groupes lisses}_{\mathrm{nilp}}(S)
\]
compatible aux changements de base, l'isomorphisme précédent étant également fonctoriel et compatible aux changements de base. Il s'ensuit donc que l'on peut définir $\underline{E}(A)_{\deleted{\unsure{S}}}$ comme $S'\handwritten{\longmapsto}\underline{E}(A)_{S'}$, avec
\[
\underline{E}(A)_{S'}=E(A')\quad,
\]
$A'$ désignant un schéma abélien qui prolonge $A$ sur $S'$, \deleted{\unsure{à}} à charge de définir \deleted{un}\inserted{un système transitif d'}isomorphismes canoniques entre les groupes obtenus pour des prolongements différents. Dans le cas où $S$ est d e carac\deleted{\unsure{é}}té\deleted{\unsure{t}}ristique nulle (donc plus question de puissances divisées !) cela est particulièrement simple; on prouve qu'il existe à isomorphisme \uline{unique} près un seul schéma en groupes lisses $E'$ sur l'épaississement $S'$ de $S$ qui prolonge $E$ (NB il suffirait même d'un nilépaississement, comme on voit par passage à la limite à partir du cas noethérien): c'est le résultat déjà signalé dans ma vieille lettre à Tate. Dans le cas général, c'est plus délicat, et peut se faire par ``la méthode de l'exponentielle'', qui sera exposée dans l\deleted{\unsure{t}}e contexte analogue des groupes de BT dans la suite du séminaire. On p\deleted{\unsure{e}}eut aussi définir $\underline{E}(A^*)$ par
\[
\underline{E}(A^*)\handwritten{=}R^1f_{A\mathrm{cris}*}(\underline{G}_{m\ A_{\mathrm{cris}}})\quad ,
\]
où $A_{\mathrm{cris}}$ \deleted{\unsure{c}}désigne le topos cristallin \deleted{\unsure{absi|abso|abs}} nilpotent absolu (relativement à $\underline{Z}$): cette méthode peut également s'adapter au cas des groupes de BT.

Quoi qu'il en soit de la \deleted{\unsure{n}}méthode utilisée pour définir $\underline{E}(A)$, on
